%%%%%%%%%%%%%%%%%%%%%%%%%%%%%%%%%%%%%%%%%%%%%%%%%%%%%%%%%%%%%%%%%%%%%%
% Lớp 11 | Đường thẳng, mặt phẳng. Quan hệ song song trong | Đường thẳng và mặt phẳng song song
% Dạng 1H4?3-7: Bài toán quỹ tích và điểm cố định
%%%%%%%%%%%%%%%%%%%%%%%%%%%%%%%%%%%%%%%%%%%%%%%%%%%%%%%%%%%%%%%%%%%%%%


% nguồn: 2026-11-15_11_1H4_TL2_Song_song (nhập 15/11/2026)
\begin{bt}%[1H4H3-7]%[ma:44230269]
% nguon: Archive/2024/SP-K10-K11-LyVanHoang_TVN-023/K11-C4-HHKG-QHSS/data/Bai-2-Dang4-Chungminhbadiemthanghangdongquy.tex #2136
%[BG - 12 New - 4in1, Lê Đạt]
	Cho tứ diện $SABC$. Trên $SA$, $SB$ và $SC$ lần lượt lấy các điểm $D$, $E$, $F$ sao cho $DE$ cắt $AB$ tại $I$, $EF$ cắt $BC$ tại $J$, $FD$ cắt $CA$ tại $K$. Chứng minh ba điểm $I$, $J$, $K$ thẳng hàng.
	\loigiai{

			\begin{center}
\begin{tikzpicture}[scale=1, font=\footnotesize, line join=round, line cap=round, >=stealth]
\coordinate (A) at (0,0);
\coordinate (C) at (5,0);
\coordinate (B) at (1.2,-1.5);
\coordinate (S) at ($(A)+(2,3)$);
\coordinate (D) at ($(S)!0.5!(A)$);
\coordinate (F) at ($(S)!2/3!(C)$);
\coordinate (E) at ($(S)!1/5!(B)$);
\fill(intersection of D--F and A--C) coordinate(K);
\fill(intersection of E--F and B--C) coordinate(J);
\fill(intersection of E--D and A--B) coordinate(I);
\draw (S)--(A)--(I)--(E)--(F)--(K)--(C)--(B)--(S)--(C) (A)--(B);
\draw[dashed](K)--(J)--(F)--(D) (I)--(J)--(C)--(A);
\foreach \x/\g in {A/190,B/-90,C/-90,S/90,I/180,K/10,J/50,F/50,E/40,D/160} \fill[black] (\x) circle (1pt) +(\g:0.3)node{$\x$};
\end{tikzpicture}	
	\end{center}
		Ta có\\
		$\bullet$ $\heva{&I=AB\cap DE\,(AB,\,DE\subset (SAB))\\&I\in AB,\, AB\subset (ABC)\\&I\in DE,\,DE\subset (DEF)}\Rightarrow I\in (ABC)\cap (DEF)$.\hfill $(1)$\\
		$\bullet$ $\heva{&K=AC\cap DF\,(AC,\,DF\subset (SAC))\\&I\in AC,\, AC\subset (ABC)\\&K\in DF,\,DF\subset (DEF)}\Rightarrow K\in (ABC)\cap (DEF)$.\hfill $(2)$\\
		$\bullet$ $\heva{&J=BC\cap EF\,(BC,\,EF\subset (SBC))\\&J\in BC,\, BC\subset (ABC)\\&J\in EF,\,EF\subset (DEF)}\Rightarrow J\in (ABC)\cap (DEF)$.\hfill $(3)$\\
		Từ $(1)$, $(2)$ và $(3)$ suy ra ba điểm $I$, $J$, $K$ thẳng hàng.
	}
\end{bt}

% nguồn: 2026-11-15_11_1H4_TL2_Song_song (nhập 15/11/2026)
\begin{bt}%[1H4H3-7]%[ma:ccd502aa]
% nguon: Archive/2024/SP-K10-K11-LyVanHoang_TVN-023/K11-C4-HHKG-QHSS/data/Bai-2-Dang4-Chungminhbadiemthanghangdongquy.tex #2137
%[BG - 12 New - 4in1, Lê Đạt]
	Cho hình chóp $S.ABCD$ có đáy là hình bình hành. Gọi $M$ là trung điểm $SB$, $N$ là điểm đối xứng của $C$ qua $M$. Gọi $P$ là điểm trên cạnh $SD$, $Q$ là giao điểm của $AP$ với $(SBC)$. Chứng minh $N$, $S$, $Q$ thẳng hàng.
	\loigiai{
		\immini
		{
			Ta có $S\in (SBC)\cap (SAD)$, $Q=AP\cap (SBC)$, $AP\subset (SAD)$ nên $Q\in (SBC)\cap(SAD)$.\\
			Suy ra $SQ=(SBC)\cap(SAD)$. Mà
			$$\heva{& BC\subset(SBC) \\ & AD\subset (SAD)\\ & BC\parallel AD}\Rightarrow SQ\parallel BC.$$
			Mặt khác $SCBN$ là hình bình hành nên $SN\parallel BC$.\\
			Vậy $N$, $S$, $Q$ thẳng hàng.
		}
		{
			\begin{tikzpicture}[scale=.7,font=\footnotesize, line join=round, line cap=round, >=stealth]
			\path
			(0,0) coordinate (A)
			(5,0) coordinate (D)
			(-2,-2) coordinate (B)
			;
			\coordinate (C) at ($(B)+(D)-(A)$);
			\coordinate (S) at ($(A)+(1,5)$);
			\coordinate (M) at ($(S)!0.5!(B)$);
			\coordinate (N) at ($(C)!2!(M)$);
			\coordinate (P) at ($(S)!1/3!(D)$);
			\coordinate (Q) at ($(A)!3/2!(P)$);
			\draw (S)--(B)--(C)--(D)--(S)--(C)--(N)--(Q)--(P);
			\draw[dashed] (S)--(A)--(B)
			(P)--(A)--(D);
			\foreach \x/\g in {A/135,B/-90,C/-60,D/0,S/90,N/120,M/200,P/90,Q/60} \fill[black](\x) circle (1pt) ($(\x)+(\g:3mm)$) node{$\x$};
			\end{tikzpicture}
		}
	}
\end{bt}

% nguồn: 2026-11-15_11_1H4_TL2_Song_song (nhập 15/11/2026)
\begin{bt}%[1H4H3-7]%[ma:1132ba8e]
% nguon: Archive/2024/SP-K10-K11-LyVanHoang_TVN-023/K11-C4-HHKG-QHSS/data/Bai-2-Dang4-Chungminhbadiemthanghangdongquy.tex #2138
%[BG - 12 New - 4in1, Lê Đạt]
	Cho hình chóp $S.ABCD$ có đáy là hình bình hành tâm $O$. Trên $SD$, $SC$ lần lượt lấy $M$ và $N$ sao cho $\dfrac{SM}{SD}=\dfrac{SN}{SC}=\dfrac{2}{3}$. Chứng minh $AN$, $BM$ và $SO$ cùng đồng quy.
	\loigiai{
		\immini
		{
			Ta có $\dfrac{SM}{SD}=\dfrac{SN}{SC}=\dfrac{2}{3}\Rightarrow MN\parallel CD$, mà $CD\parallel AB$ nên $MN\parallel AB$ và $M$, $N$, $B$, $A$ đồng phẳng.\\
			Gọi $I$là giao điểm của $AN$ và $BM$ suy ra $I\in (SAC)\cap (SBD)$.\\
			Mà $SO=(SAC)\cap (SBD)$ nên $SO$ đi qua $I$.\\
			Vậy $AN$, $BM$ và $SO$ cùng đồng quy tại $I$.
		}
		{
			\begin{tikzpicture}[scale=.7,font=\footnotesize, line join=round, line cap=round, >=stealth]
			\path
			(0,0) coordinate (A)
			(5,0) coordinate (D)
			(-2,-2) coordinate (B)
			;				
			\coordinate (C) at ($(B)+(D)-(A)$);
			\coordinate (S) at ($(A)+(1,5)$);
			\coordinate (M) at ($(S)!2/3!(D)$);
			\coordinate (N) at ($(C)!1/3!(S)$);
			\coordinate (O) at ($(A)!.5!(C)$);
			\coordinate (I) at ($(A)!3/5!(N)$);
			\draw (S)--(B)--(C)--(D)--(S)--(C)
			(M)--(N);
			\draw[dashed] (O)--(S)--(A)--(C)
			(D)--(B)--(A)--(D)
			(B)--(M)
			(A)--(N);
			\foreach \x/\g in {A/180,B/-90,C/-60,D/0,S/90,N/-10,M/60,I/120,O/-90} \fill[black](\x) circle (1pt) ($(\x)+(\g:3mm)$) node{$\x$};
			\end{tikzpicture}
		}
	}
\end{bt}

% nguồn: 2026-11-15_11_1H4_TL2_Song_song (nhập 15/11/2026)
\begin{bt}%[1H4H3-7]%[ma:5dfff45d]
% nguon: Archive/2024/SP-K10-K11-LyVanHoang_TVN-023/K11-C4-HHKG-QHSS/data/Bai-2-Dang4-Chungminhbadiemthanghangdongquy.tex #2142
%[BG - 12 New - 4in1, Lê Đạt]
	Cho hình chóp $S.ABCD$ có đáy là hình thang. Gọi $E$, $F$, $I$, $J$ lần lượt là trung điểm $SC$, $SD$, $SA$, $SB$. Gọi $O$ là giao điểm của $AC$ và $BD$, $K$ là giao điểm của $IE$ và $FI$. Chứng minh $S$, $K$, $O$ thẳng hàng.
	\loigiai{
		\immini
		{
			Giả sử $AB\parallel CD$. Mà theo tính chất đường trung bình thì $EF\parallel CD$, $IJ\parallel AB$ nên $IJ\parallel FE$ và $E$, $F$, $I$, $J$ đồng phẳng.\\
			Xét $3$ mặt phẳng $(SAC)$, $(SBD)$ và $(EFIJ)$ có  $\heva{& (SAC)\cap (SBD)=SO \\ & (SAC)\cap (EFIJ)=IE\\ & (SBD)\cap (EFIJ)=FJ.}$\\
			Mà $IE$ cắt $FJ$ nên $SO$, $IE$ và $FJ$ đồng quy tại $O$ hay $S$, $K$, $O$ thẳng hàng.
		}
		{
			\begin{tikzpicture}[scale=1, font=\footnotesize,line join=round, line cap=round, >=stealth]
			\path
			(0,0) coordinate (A)
			(-1.3,-1.2) coordinate (D)
			(2.5,-1.2) coordinate (C)
			;
			\coordinate (B) at ($(A)+(C)-(D)$);
			\coordinate (S) at ($(A)+(0.5,3)$);
			\coordinate (I) at ($(S)!1/2!(A)$);
			\coordinate (J) at ($(S)!1/2!(B)$);
			\coordinate (E) at ($(S)!1/2!(C)$);
			\coordinate (F) at ($(S)!1/2!(D)$);
			\coordinate (K) at ($(I)!0.5!(E)$);
			\coordinate (O) at ($(A)!0.5!(C)$);
			\draw (S)--(B)--(C)--(D)--cycle
			(S)--(C) (J)--(E)--(F);
			\draw[dashed] (O)--(S)--(A)--(B)--(D)--(A)--(C)
			(E)--(I)--(J)--(F)--(I); 
			\foreach \p/\q in {A/160,B/0,C/-90,D/-90,S/90,I/45,J/45,K/45, E/-30,F/180,O/-90}
			\fill[black] (\p) circle (1.0pt) ($(\p)+(\q:3.5mm)$) node{$\p$};
			\end{tikzpicture}}
	}
\end{bt}

% nguồn: 2026-11-15_11_1H4_TL2_Song_song (nhập 15/11/2026)
\begin{bt}%[1H4H3-7]%[ma:51d802f0]
% nguon: Archive/2024/SP-K10-K11-LyVanHoang_TVN-023/K11-C4-HHKG-QHSS/data/Bai-2-Dang4-Chungminhbadiemthanghangdongquy.tex #2143
%[BG - 12 New - 4in1, Lê Đạt]
	Cho hình chóp $S.ABC$. Gọi $K,N$ lần lượt là trung điểm của $SA,BC$. Điểm $M$ thuộc cạnh $SC$ sao cho $SM =\dfrac{2}{3}SC$.
	\begin{enumerate}
		\item Tìm giao tuyến của hai mặt phẳng $(KMN)$ và $(SBC)$.
		\item Tìm giao điểm của đường thẳng $AC$ và mặt phẳng $(KMN)$. 
		\item Gọi $H$ là giao điểm của đường thẳng $AB$ và mặt phẳng $(KMN)$. Chứng minh rằng ba đường thẳng $KN$, $MN$, $SB$ đồng quy.
	\end{enumerate}
	\loigiai{
				\begin{center}
		\begin{tikzpicture}[scale=1, font=\footnotesize, line join=round, line cap=round, >=stealth]
		\coordinate (A) at (0,0);
		\coordinate (C) at (5,0);
		\coordinate (B) at (1.2,-1.5);
		\coordinate (S) at ($(A)+(2,3)$);
		\coordinate (K) at ($(S)!0.5!(A)$);
		\coordinate (M) at ($(S)!2/3!(C)$);
		\coordinate (H) at ($(A)!3/4!(B)$);
		\fill(intersection of K--H and S--B) coordinate(E);
		\fill(intersection of E--M and B--C) coordinate(N);
		\fill(intersection of K--M and A--C) coordinate(D);
		\draw (S)--(A)--(B)--(C)--(S)--(E)--(M)--(D)--(C) (K)--(E);
		\draw[dashed](N)--(K)--(M) (A)--(C);
		\foreach \x/\g in {A/190,B/-50,C/-90,S/90,K/180,E/0,M/50,N/-50,H/-140,D/90} \fill[black] (\x) circle (1pt) +(\g:0.3)node{$\x$};
		\end{tikzpicture}	
	\end{center}
		\begin{enumerate}
			\item Ta có $(KMN) \cap (SBC) = MN$.
			\item Trong mặt phẳng $(SAC)$, gọi $D$ là giao điểm của $AC$ và $KM$, suy ra $AC \cap (KMN) = \{ D\}$.
			\item Trong mặt phẳng $(SBC)$, gọi $E$ là giao điểm của $MN$ và $SB$, do đó trong mặt phẳng $(SAB)$, $H$ là giao điểm của $AB$ và $KE$. Dễ thấy ba đường thẳng $KN$, $MN$, $SB$ đồng quy tại $E$.
		\end{enumerate}
	}
\end{bt}

% nguồn: 2026-11-15_11_1H4_TL2_Song_song (nhập 15/11/2026)
\begin{bt}%[1H4H3-7]%[ma:4b7ef85d]
% nguon: Archive/2024/SP-K10-K11-LyVanHoang_TVN-023/K11-C4-HHKG-QHSS/data/Bai-2-Dang4-Chungminhbadiemthanghangdongquy.tex #2144
%[BG - 12 New - 4in1, Lê Đạt]
	Cho hình chóp $S.ABCD$ đáy là hình thang đáy lớn $AD$, gọi $K$ là giao điểm của $AB$ và $CD$, $M$ là điểm nằm trên $SD$, $P$ là giao điểm của $KM$ và $SC$. Gọi $d$ là giao tuyến của $(SAD)$ và $(SBC)$. Chứng minh $AM$, $BP$ và $d$ đồng quy.
	\loigiai{
		\immini
		{
			Ta có
			$$\heva{& AD\subset(SAD) \\ & BC\subset (SBC)\\ & AD\parallel BC\\ & S\in (SAD)\cap (SBC)}\Rightarrow (SAD)\cap (SBC)=Sx\parallel AD\parallel BC.$$
			Xét $3$ mặt phẳng $(SAD)$, $(SBC)$, $(AMK)$ có
			$\heva{& (SAD)\cap (SBC)=d \\ & (SAD)\cap (AMK)=AM\\ & (AMK)\cap (SBC)=BP.}$\\
			Mà $AM$ cắt $Sx$ hay $AM$ cắt $d$, vậy $AM$, $BP$ và $d$ đồng quy.
		}
		{
			\begin{tikzpicture}[scale=.7,font=\footnotesize, line join=round, line cap=round, >=stealth]
			\path
			(0,0) coordinate (A)
			(5,0) coordinate (D)
			(1,-2) coordinate (B)
			(3.5,-2) coordinate (C)
			;				
			\coordinate (O) at ($(A)!2/3!(C)$);
			\coordinate (S) at ($(O)+(-.5,6)$);
			\coordinate (M) at ($(S)!0.5!(D)$);
			\coordinate (K) at ($(A)!2!(B)$);
			\coordinate (P) at ($(S)!2/3!(C)$);
			\coordinate (I) at ($(A)!2!(M)$);
			\path
			(intersection of B--P and S--D) coordinate (x)
			;
			\draw (S)--(A)--(K)--(D)--(x) (S)--(I)--(B)--(S)--(C)
			(K)--(P)
			(B)--(I);
			\draw[dashed] (I)--(A)--(D)
			(B)--(C)
			(M)--(P) (x)--(S);
			\foreach \x/\g in {A/180,B/-120,C/-60,D/0,S/90,K/-90,P/20,I/90,M/70} \fill[black](\x) circle (1pt) ($(\x)+(\g:4mm)$) node{$\x$};
			\end{tikzpicture}
		}
	}
\end{bt}

% nguồn: 2026-11-15_11_1H4_TL2_Song_song (nhập 15/11/2026)
\begin{bt}%[1H4H3-7]%[ma:49a795f5]
% nguon: Archive/2024/SP-K10-K11-LyVanHoang_TVN-023/K11-C4-HHKG-QHSS/data/Bai-2-Dang4-Chungminhbadiemthanghangdongquy.tex #2145
%[BG - 12 New - 4in1, Lê Đạt]
	Cho hình chóp $S.ABCD$ có đáy $ABCD$ là tứ giác có các cạnh đối không song song.
	\begin{enumerate}
		\item Tìm giao tuyến của $(SAC)$ và $(SBD)$.
		\item Gọi $M$ là điểm trên cạnh $SC$ ($M$ không trùng với $S$ và $C$). Tìm giao điểm của $AM$ và $(SBD)$.
		\item Gọi $H$ là giao điểm của $SD$ và $(ABM)$, $I$ là giao điểm của $BM$ và $(SAD)$. Chứng minh rằng $A$, $I$, $H$ thẳng hàng.
	\end{enumerate}
	\loigiai{
	\begin{center}
		\begin{tikzpicture}[scale=1,font=\footnotesize, line join=round, line cap=round, >=stealth]
		\coordinate (A) at (0,0);
		\coordinate (D) at (5,0);
		\coordinate (B) at (1.5,-3);
		\coordinate (C) at (4,-2);
		\coordinate (S) at ($(A)+(1,4)$);
		\coordinate (M) at ($(S)!0.6!(C)$);
		\fill(intersection of A--C and B--D) coordinate(O);
		\fill(intersection of S--O and A--M) coordinate(E);
		\fill(intersection of B--E and S--D) coordinate(H);
		\fill(intersection of A--D and B--C) coordinate(F);
		\fill(intersection of A--H and S--F) coordinate(I);
		\foreach \x/\g in {A/160,B/-100,C/-90,D/20,S/90,O/190,M/0,H/90,E/160,I/80,F/90} \fill[black](\x) circle (1.5pt) ($(\x)+(\g:3mm)$) node{$\x$};
		\draw (S)--(A)--(B)--(C)--(S)--(B)--(I) (S)--(F)--(C);
		\draw[dashed] (F)--(A)--(C)--(D)--(B)--(H)--(A)--(M) (D)--(S)--(O)
		(H)--(I);
		\end{tikzpicture}
	\end{center}
		\begin{enumerate}
			\item Tìm giao tuyến của $(SAC)$ và $(SBD)$.\\
			Trong mặt phẳng $(ABCD)$ gọi $O$ là giao điểm của $AC$ và $BD$. Xét hai mặt phẳng $(SAC)$ và $(SBD)$. Ta có\\
			$S$ là một điểm chung của hai mặt phẳng.\\
			$O \in AC \Leftrightarrow O \in (SAC)$, $O \in BD \Leftrightarrow O \in (SBD)$ nên $O$ là điểm chung thứ hai.\\
			Mà $S \notin (ABCD)$ nên $S \neq O$.\\
			Vậy $(SAC) \cap (SBD)=SO$.
			\item Gọi $M$ là điểm trên cạnh $SC$ ($M$ không trùng với $S$ và $C$). Tìm giao điểm của $AM$ và $(SBD)$.\\
			Do $M \in SC$ nên $AM \subset (SAC)$. Trong mặt phẳng $(SAC)$, gọi $E$ là giao điểm của hai đường thẳng $SO$ và $AM$.\\
			Ta có $E \in SO,\ SO \subset (SBD) \Rightarrow E \in (SBD)$ và $E \in AM$.\\
			Vậy $E$ là giao điểm của $AM$ và mặt phẳng $(SBD)$.
			\item Gọi $H$ là giao điểm của $SD$ và $(ABM)$, $I$ là giao điểm của $BM$ và $(SAD)$. Chứng minh rằng $A$, $I$, $H$ thẳng hàng.\\
			Trong mặt phẳng $(ABCD)$, do $AD$ và $BC$ không song song nên chúng cắt nhau giả sử tại $F$. Ta có $SF \subset(SAD)$ và $SF \subset (SBC)$.\\
			Trong mặt phẳng $(SBC)$, gọi $J$ là giao điểm của $BM$ và $SF$, khi đó do $SF \subset (SAD)$ nên $J$ là điểm chung của hai mặt phẳng $(ABM)$ và $(SAD)$.\\
			Mặt khác mặt phẳng $(ABM)$ và $(SAD)$ phân biệt, có $A$ là điểm chung, $A \not\equiv J$.\\
			Giao tuyến của $(ABM)$ và $(SAD)$ là đường thẳng $\Delta$ qua hai điểm $A, J$.\\
			Theo bài ra $H$ là giao điểm của $SD$ và $(ABM)$ nên $H \in (ABM)$, $H \in SD \Rightarrow H \in (SAD)$. Do đó $H$ là một điểm chung của $(ABM)$ và $(SAD)$. Suy ra $H \in \Delta$.\\
			Lại có $\{I\}=BM \cap (SAD)$ nên $I \equiv J$. Suy ra $I \in \Delta$.\\
			Vậy $A$, $H$, $I$ thẳng hàng.
		\end{enumerate}			
	}
\end{bt}
